\documentclass[aspectratio=169]{beamer}
\usepackage[utf8]{inputenc}\usepackage{amsmath}\usepackage{tcolorbox}
\definecolor{UNIBBlue}{RGB}{0, 51, 102}\setbeamercolor{frametitle}{bg=UNIBBlue,fg=white}
\title{Optimisasi untuk Teknik Elektro}\subtitle{Minggu 7: Optimisasi Cembung Terkendala (KKT)}
\date{Semester Ganjil 2026}\begin{document}\begin{frame}\titlepage\end{frame}
\begin{frame}{Kondisi Karush-Kuhn-Tucker (KKT)}
Generalisasi dari metode Pengali Lagrange untuk menangani kendala pertidaksamaan.
Syarat perlu (dan cukup untuk masalah konveks) agar suatu solusi menjadi optimal.
\begin{itemize}
\item Stationarity: Gradien Lagrangia bernilai 0.
\item Primal Feasibility: Memenuhi batasan asli.
\item Dual Feasibility: Pengali Lagrange $\ge 0$ untuk pertidaksamaan.
\item Complementary Slackness: $\lambda_i \cdot g_i(x) = 0$.
\end{itemize}
\end{frame}\end{document}